\chapter{Deployment and Go-Live}\label{ch:deploy}

\section{Deployment Environment}

Infra Monitor is containerised with Docker Compose and deployable on any
Linux host (Ubuntu 22.04 LTS recommended) with Docker Engine 24+ installed.
For this academic submission, the \textbf{Django backend has been deployed to
Render} (a cloud PaaS), making the REST API and WebSocket endpoint publicly
accessible.  VictoriaMetrics (the time-series database cluster) was not deployed
to a public host because the VictoriaMetrics cluster mode requires persistent
volumes and a paid tier on most PaaS providers; it continues to run locally
via Docker Compose during development and testing.

The full local stack consists of nine services defined in \texttt{docker-compose.yml}:

\begin{table}[H]
\centering
\caption{Docker Compose service inventory}\label{tab:services}
\begin{tabularx}{\textwidth}{@{}l l X l@{}}
\toprule
\textbf{Service} & \textbf{Image / Build} & \textbf{Role} & \textbf{Port}\\
\midrule
backend        & \texttt{./backend} & Django+Daphne ASGI API & 7000:8000\\
ml\_service    & \texttt{./model}  & FastAPI ML inference microservice & 7001:80\\
celery\_worker & \texttt{./backend} & Celery async task executor & --\\
celery\_beat   & \texttt{./backend} & Periodic task scheduler & --\\
redis          & redis:8-alpine     & Celery broker + WS channel layer & --\\
postgres       & postgres:18        & Primary relational database & 5433:5432\\
vmstorage      & victoriametrics/vmstorage:v1.150.0-cluster & Time-series storage & --\\
vminsert       & victoriametrics/vminsert:v1.150.0-cluster  & Metric write endpoint & --\\
vmselect       & victoriametrics/vmselect:v1.150.0-cluster  & Metric query endpoint & 8481\\
\bottomrule
\end{tabularx}
\end{table}

\textbf{Data volumes:}
\begin{itemize}
  \item \texttt{postgres-data} -- PostgreSQL data directory
  \item \texttt{victoriametrics-data} -- VictoriaMetrics storage
  \item \texttt{ml\_artifacts} -- Trained model joblib files
\end{itemize}

\textbf{Health checks:} All stateful services define Docker health checks
(PostgreSQL: \texttt{pg\_isready}; Redis: \texttt{redis-cli ping};
VictoriaMetrics: HTTP \texttt{/health}; ML service: HTTP \texttt{/health}).
The backend service depends on all three external services being healthy
before starting (\texttt{condition: service\_healthy}).

\textbf{Environment configuration:}
\begin{itemize}
  \item \texttt{backend/.env}: \texttt{SECRET\_KEY}, \texttt{DATABASE\_URL},
        \texttt{GEMINI\_API\_KEY}, \texttt{ML\_SERVICE\_TOKEN},
        \texttt{FIREBASE\_CREDENTIALS}, \texttt{ALLOWED\_HOSTS},
        \texttt{CELERY\_BROKER\_URL}
  \item \texttt{model/.env}: \texttt{ML\_SERVICE\_TOKEN},
        \texttt{DJANGO\_INTERNAL\_URL}
\end{itemize}

\textbf{SSL/TLS:} In production, a reverse proxy (Nginx or Caddy) is placed
in front of the backend container to terminate HTTPS and proxy WebSocket
upgrades.  Let's Encrypt or an internal CA provides the certificate.

\textbf{Deployment diagram:} See \Cref{fig:deploy_diagram} in Appendix~E.

\section{Deployment Process and CI/CD Implementation}

\textbf{Step-by-step deployment procedure:}

\begin{enumerate}
  \item \textbf{Provision host:} Launch an Ubuntu 22.04 LTS VM or bare-metal
        server (minimum 4~GB RAM, 20~GB SSD).
  \item \textbf{Install Docker:}
\begin{lstlisting}[language=bash,caption={Docker installation (Ubuntu)},label={lst:docker}]
sudo apt update
sudo apt install -y docker.io docker-compose-plugin
sudo systemctl enable --now docker
sudo usermod -aG docker $USER
\end{lstlisting}
  \item \textbf{Clone repository:}
\begin{lstlisting}[language=bash,caption={Clone and configure},label={lst:clone}]
git clone https://github.com/[ASSUMED: repo-url] infra-monitor
cd infra-monitor
cp backend/.env.example backend/.env
cp model/.env.example model/.env
# Edit both .env files with actual secrets
\end{lstlisting}
  \item \textbf{Build and start services:}
\begin{lstlisting}[language=bash,caption={Start the full stack},label={lst:compose}]
docker compose build
docker compose up -d
\end{lstlisting}
  \item \textbf{Run migrations and create superuser:}
\begin{lstlisting}[language=bash,caption={Database initialisation},label={lst:migrate}]
docker compose exec backend python manage.py migrate
docker compose exec backend python manage.py createsuperuser
\end{lstlisting}
  \item \textbf{Train initial ML model:}
\begin{lstlisting}[language=bash,caption={Initial ML model training},label={lst:train}]
# Trigger training via Celery or admin API
docker compose exec backend python manage.py trigger_ml_training
\end{lstlisting}
  \item \textbf{Configure reverse proxy:} Point Nginx to port 7000 (backend)
        with WebSocket upgrade support (\texttt{Upgrade: websocket} header
        pass-through).
  \item \textbf{Verify health:}
\begin{lstlisting}[language=bash,caption={Health verification},label={lst:health}]
curl http://localhost:7000/api/health/
# Expected: {"status": "ok"}
curl http://localhost:7001/health
# Expected: {"status": "ok"}
\end{lstlisting}
\end{enumerate}

\section{Go-Live Strategy and Post-Deployment Monitoring}

\textbf{Academic deployment summary:} The Django API backend has been deployed
to Render for the academic submission.  The Docker Compose configuration
(\texttt{docker-compose.yml}) is verified and all services start correctly
locally with \texttt{docker compose up}.  VictoriaMetrics is not publicly hosted
at this stage; it runs locally for development and testing purposes, as hosting
the cluster mode on a public platform would require a paid persistent-storage tier.

\textbf{Docker configuration status:} All Compose services have health checks
(PostgreSQL: \texttt{pg\_isready}; Redis: \texttt{redis-cli ping};
VictoriaMetrics: HTTP \texttt{/health}; ML service: HTTP \texttt{/health}).
Service startup ordering is enforced with \texttt{condition: service\_healthy}.
Environment variables are loaded from \texttt{.env} files; no secrets are
committed to the repository.

\textbf{Rollout strategy:} For a full production deployment, a \textbf{phased
rollout} is recommended: deploy monitoring-only features first (server
enrolment, dashboards), then enable alerts/incidents, then enable AI features.

\textbf{Post-deployment monitoring (planned):}
\begin{itemize}
  \item \textbf{Application health:} \texttt{GET /api/health/} polled every
        30~s by an external uptime monitor.
  \item \textbf{Container health:} \texttt{docker compose ps} and Docker
        restart policies (\texttt{restart: unless-stopped}) ensure automatic
        recovery from crashes.
  \item \textbf{Log monitoring:} \texttt{docker compose logs -f backend}
        for real-time error tracking.
  \item \textbf{Rollback plan:} Tag each release (\texttt{git tag v1.x.x});
        restore previous image with \texttt{docker compose pull} and
        \texttt{docker compose up -d}.
\end{itemize}

